\section*{Status, attribution, and evidence boundary}
This is an AI-generated, internally audited research draft for expert
examination, not a submitted or peer-reviewed article. The project initiator
reports providing time, a subscription, and local computing resources, but
no personal mathematical contribution or capacity to verify the proofs.
No accountable scholarly author has yet been established. An empty author
field is not an attribution decision, and an AI system is not listed as an
author. Rights remain reserved; author responsibility, declarations,
literature priority, and independent expert review remain publication gates.

This compact paper is one of three reading editions of the same research
project, not a third independent publication. The \emph{Selected Results}
edition retains the former extended compact companion. The \emph{Full
Research Report} retains the comprehensive technical account. All three
use the evidence cutoff of 30 September 2026 at 12:54 UTC. The present
selection is deliberately narrower: it develops the direct-product
theorems, a binary-congruence construction, and small-order examples.
It does not append every later construction or search campaign.

Statements with written proofs are distinguished from those labelled
\emph{Computational Theorem}. The latter depend on exact finite evidence;
the order-$8$ census also depends on the completeness of the cited external
enumeration. Fresh checks of displayed tables do not rerun an entire census.
Internal proof review, independent implementations, and machine proof
checking are not independent human refereeing.

The order-$12$ witness refutes the former power-of-two-only conjecture.
Order $10$ is excluded in the full report by a separate mathematical
reduction and exact computation; that result is background here, not a
proof supplied by this compact paper. None of the three editions resolves
unrestricted order $18$, proves a complete higher-order main-class census,
or establishes literature priority. Changes in selection and presentation
do not change the evidence status of the underlying results.
